
All that remains to prove \thmref{bayes_clt_main} is to show that the residual is
$O_p(1)$ when summed.  The form of $\resid{\phi}(\xvec_s)$ is awkward, but it is
designed specifically to commute with suprema over sums, allowing us to apply
ULLNs to the remainder of our BCLT, as stated formally in
\lemref{integral_commuting_sums,taylor_commuting_sums} below.

Observe that residual expressions of the form $\sup_{\theta \in R_1 \bigcup R_2}
\norm{\phi(\theta, \x)}_2$ would not be as useful, since
%
\begin{align*}
%
\meann \sup_{\theta \in R_1 \bigcup R_2} \norm{\phi(\theta, \x_n)}_2 \ge
\sup_{\theta \in R_1 \bigcup R_2} \meann \norm{\phi(\theta, \x_n)}_2,
\end{align*}
%
and ULLN bound on the latter does not imply a bound on the former, since the
supremum may depend on $\x_n$ in a way that causes the sample mean to diverge.

%%%%%%%%%%%%%%%%%%%%%%%%%%%%%%%%%%%%%%%%%%%%%%%%%%%%%%%%%%%%%%%%%%%%%%%
%%%%%%%%%%%%%%%%%%%%%%%%%%%%%%%%%%%%%%%%%%%%%%%%%%%%%%%%%%%%%%%%%%%%%%%
%%%%%%%%%%%%%%%%%%%%%%%%%%%%%%%%%%%%%%%%%%%%%%%%%%%%%%%%%%%%%%%%%%%%%%%

\begin{lem}\lemlabel{integral_commuting_sums}
%
The Gaussian integrals of \defref{bclt_resid} commutes with sums.
%
For a scalar--valued function $\phi(\theta, \xvec_s)$ and an index set
$\indexset$,
%
\begin{align*}
%
\means \expect{\normresid{\infoev / 2}{R_1 \bigcup R_2}(\tau)}
              {\phi(\theta, \xvec_s)}^2
% \means \resid{2}(\xvec_s)
            \le
\sup_{\theta \in R_1 \bigcup R_2} \means \phi(\theta, \xvec_s)^2.
%
\end{align*}
%
\begin{proof}
%
\begin{align*}
%
\means \expect{\normresid{\infoev / 2}{R_1 \bigcup R_2}(\tau)}
              {\phi(\theta, \xvec_s)}^2
% \means \resid{2}(\xvec_s)
% ={}&
% \means \expect{\normresid{\infoev / 2}{R_1 \bigcup R_2}(\tau)}
%             {\phi(\theta, \xvec_s)^2}
\le{}&
\expect{\normresid{\infoev / 2}{R_1 \bigcup R_2}(\tau)}
              {\means \phi(\theta, \xvec_s)^2}
\\\le{}&
\sup_{\tau \in R_1 \bigcup R_2}
    \means \phi(\thetahat + \tau / \sqrt{N}, \xvec_s)^2
\\={}&
\sup_{\theta \in R_1 \bigcup R_2} \means \phi(\theta, \xvec_s)^2.
%
\end{align*}
%
\end{proof}
%
\end{lem}
%%%%%%%%%%%%%%%%%%%%%%%%%%%%%%%%%%%%%%%%%%%%%%%%%%%%%%%%%%%%%%%%%%%%%%%


%%%%%%%%%%%%%%%%%%%%%%%%%%%%%%%%%%%%%%%%%%%%%%%%%%%%%%%%%%%%%%%%%%%%%%%
%%%%%%%%%%%%%%%%%%%%%%%%%%%%%%%%%%%%%%%%%%%%%%%%%%%%%%%%%%%%%%%%%%%%%%%
%%%%%%%%%%%%%%%%%%%%%%%%%%%%%%%%%%%%%%%%%%%%%%%%%%%%%%%%%%%%%%%%%%%%%%%


\begin{lem}\lemlabel{taylor_commuting_sums}
%
The Taylor series residual of \defref{bclt_resid} commutes with sums.
%
When $\theta \in R_1 \bigcup R_2$, for a k-th order continuously differentiable
function $\phi(\theta, \xvec_s)$ and an index set $\indexset$,
%
\begin{align*}
%
\means \norm{\tresid{k}{\phi(\cdot, \xvec_s), \theta, \thetahat}}_2^2 \le
\sup_{\theta \in R_1 \bigcup R_2}
    \means \norm{\phi_{(k)}(\theta, \xvec_s)}_2^2.
%
\end{align*}
%
\begin{proof}
%
Recall from \defref{taylor_residual} that
%
\begin{align*}
%
\MoveEqLeft
\means \norm{\tresid{k}{\phi(\cdot, \xvec_s), \theta, \thetahat}}_2^2
=\\{}&
\means \norm{
    \frac{1}{(k-1)!} \int_0^1 (1 - t)^{k-1}
    \phi_{(k)}(\thetahat + t (\theta - \thetahat), \xvec_s) dt
}_2^2
\le\\&
\means \frac{1}{(k-1)!} \int_0^1
\norm{\phi_{(k)}(\thetahat + t (\theta - \thetahat), \xvec_s)}_2^2 dt
=\\&
\frac{1}{(k-1)!} \int_0^1
\means \norm{\phi_{(k)}(\thetahat + t (\theta - \thetahat), \xvec_s)}_2^2 dt
\le\\&
\frac{1}{(k-1)!}
\sup_{\theta \in R_1 \bigcup R_2}
    \means\norm{\phi_{(k)}(\theta, \xvec_s)}_2^2.
%
\end{align*}
%
The result follows because $k \ge 1$.
%
\end{proof}
%
\end{lem}

%%%%%%%%%%%%%%%%%%%%%%%%%%%%%%%%%%%%%%%%%%%%%%%%%%%%%%%%%%%%%%%%%%%%%%%


%%%%%%%%%%%%%%%%%%%%%%%%%%%%%%%%%%%%%%%%%%%%%%%%%%%%%%%%%%%%%%%%%%%%%%%
%%%%%%%%%%%%%%%%%%%%%%%%%%%%%%%%%%%%%%%%%%%%%%%%%%%%%%%%%%%%%%%%%%%%%%%
%%%%%%%%%%%%%%%%%%%%%%%%%%%%%%%%%%%%%%%%%%%%%%%%%%%%%%%%%%%%%%%%%%%%%%%

%
\begin{lem}\lemlabel{resid_op1}
%
Under \assuref{core_bclt_assu}, $\meann\resid{\phi}(\x_n)^2 = O_p(1)$.
%
\begin{proof}
%
% By \lemref{regular_event} (in particular \lemref{regular_event} \eventref{ulln}),
% together with

% TODO: it would be good to just apply the lemmas once to get a big
% expression in terms of suprema over \thetaball{\delta}.

We may take $\delta_2$ to be the smaller of the values required for
\defref{bclt_okay} \itemref{grad_op1} and
the region $R_2$ \lemref{region_r2} to
hold.  Then, by \defref{bclt_okay_index} \itemref{grad_op1_index},
for $k=0,1,2,3$,
%
\begin{align*}
%
\sup_{\theta \in R_1 \bigcup R_2}
    \meann \norm{\phigrad{k}(\theta, \x_n)}_2^2 = O_p(1),
%
\end{align*}
%

It follows immediately that $\meann
\norm{\phigrad{k}(\thetahat, \x_n)}_2 = O_p(1)$ and $\meann
\norm{\phigrad{k}(\thetahat, \x_n)}_2^2 = O_p(1)$ for $k=0, 1, 2, 3$. From this
and Cauchy-Schwarz it follows that $\meann \resid{0}(\x_n)^2 = O_p(1)$.

Applying \lemref{integral_commuting_sums} with $\indexset = \{1, \ldots, N \}$ and
$\phi(\theta, \xvec_s) = \norm{\phi(\theta, \x_n)}_2$ gives
$\meann \resid{2}(\x_n) = O_p(1)$ and $\meann \resid{2}(\x_n)^2 = O_p(1)$.

Similarly, $\meann \resid{1}(\x_n)^2 = O_p(1)$ follows from
\lemref{integral_commuting_sums} and then \lemref{taylor_commuting_sums}.

The terms $\meann \resid{3}(\x_n)$ and $\meann \resid{3}(\x_n)^2$ are controlled
by assumption in \defref{bclt_okay} \itemref{prior_op1}.

Finally, the fact that $\meann \resid{\phi}(\x_n)^2 = O_p(1)$ follows
by Cauchy-Schwarz, since $\resid{\phi}(\x_n)$ is the sum of products
of terms, the squares of each of which have $O_p(1)$ sum.
%
\end{proof}
%
\end{lem}
%
%%%%%%%%%%%%%%%%%%%%%%%%%%%%%%%%%%%%%%%%%%%%%%%%%%%%%%%%%%%%%%%%
